\section{The model and the limiting law}\label{sec:model}\label{mod:section}

\subsection{Domains, weights and curves}
Fix a bounded simply connected domain $\Omega\subset\C$, an interior
point $z_*$, distinct boundary prime ends $a,b$, and a real mass
$m\ne0$. Throughout, boundary statements for a simply connected domain
refer to its prime-end compactification; see \cite{Pommerenke}
for this standard conformal background. No local connectivity or
accessibility hypothesis is imposed on its physical boundary.

For each mesh $\delta$, let $\Omega^\delta$ be a finite simply connected
polygonal domain made from rhombi of an isoradial lattice. Every primal face has
circumradius $\delta$; the dual vertices are the circumcentres. A primal
edge and its dual form the diagonals of a rhombus. Its half-angle at a
primal vertex, denoted by $\theta_e$, satisfies
\begin{equation}\label{mod:angles}
 \eta\le\theta_e\le\frac\pi2-\eta
\end{equation}
for one $\eta>0$ independent of the mesh. Boundary half-rhombi and
boundary occurrences of vertices are interpreted with the usual wired
edge convention, described further in Section~\ref{sec:geometry}.
The lattices themselves may vary with $\delta$.

Let $\chi_\delta:\Omega^\delta\to\D$ and
$\chi:\Omega\to\D$ be normalized at $z_*$ by
$\chi_\delta(z_*)=\chi(z_*)=0$ and positive real derivative. For all
sufficiently small meshes the basepoint belongs to $\Omega^\delta$.
The marked convergence assumed here is precisely kernel convergence
from $z_*$, together with
\begin{equation}\label{mod:marks}
 \chi_\delta(a^\delta)\longrightarrow\chi(a),\qquad
 \chi_\delta(b^\delta)\longrightarrow\chi(b).
\end{equation}
In particular, there is no common diameter or area bound on the
approximating domains.

For $0<k<1$ put
\[
 U_k(t)=\int_0^t\frac{\dd s}{\sqrt{1-k^2\sin^2s}},\qquad
 K(k)=U_k(\pi/2).
\]
Choose any nome sequence $r_\delta\in(0,1)$ such that
$r_\delta/\delta\to |m|/2$, and define $k_\delta\in(0,1)$ by
\[
 r_\delta=\exp\left\{-\pi\frac{K(\sqrt{1-k_\delta^2})}
 {K(k_\delta)}\right\}.
\]
Let $A_\delta(\theta)\in(0,\pi/2)$ be determined by
$U_{k_\delta}(A_\delta(\theta))=2K(k_\delta)\theta/\pi$ and set
\begin{equation}\label{mod:thermal-angle}
 \widehat\theta_e=
 \begin{cases}
 A_\delta(\theta_e),&m>0,\\
 \pi/2-A_\delta(\pi/2-\theta_e),&m<0.
 \end{cases}
\end{equation}

Wire the counterclockwise boundary arc $(b^\delta a^\delta)$ by
primal-open edges and $(a^\delta b^\delta)$ by dual-open edges. For a
set $E$ of interior primal-open edges, put a dual-open edge across each
edge outside $E$. The smoothed cluster boundaries comprise one
oriented interface $\gamma_\delta:a^\delta\to b^\delta$ and
$\ell(E)$ closed medial loops. The probability of $E$ is
\begin{equation}\label{mod:law}
 \Prob_{\delta,m}(E)=\frac1{Z_{\delta,m}}(\sqrt2)^{\ell(E)}
 \prod_{e\in E}\sin\frac{\widehat\theta_e}{2}
 \prod_{e\notin E}\sin\left(\frac\pi4-
                                    \frac{\widehat\theta_e}{2}\right).
\end{equation}
Here $Z_{\delta,m}$ normalizes the sum over all subsets of interior
primal edges.
The cluster weight is $2$. The symbol $r_\delta$ denotes the elliptic
nome and has no relation to the cluster weight. The negative branch
in Equation~\eqref{mod:thermal-angle} is exactly the planar dual of
the positive branch.

All curve laws are compared in the same space: continuous oriented
curves in $\overline\D$ modulo increasing reparametrization, equipped
with the uniform curve metric
\begin{equation}\label{mod:curve-metric}
 d_{\mathrm{curv}}(\alpha,\beta)
 =\inf_{\rho}\sup_{0\le s\le1}|\alpha(s)-\beta(\rho(s))|,
\end{equation}
where the infimum is over increasing homeomorphisms of $[0,1]$;
curves at distance zero are identified. Equivalently, one may allow
continuous nondecreasing surjections and parametrize both curves.
The curve representing the discrete model is
$\chi_\delta\circ\gamma_\delta$. Thus a physically long piece in a
vanishing disc cap has vanishing curve diameter in the topology used
here.

\begin{theorem}\label{thm:main}
For every $(\Omega,a,b,m)$ as above, there is a curve law
$\mu_{\Omega,a,b,m}$ such that, for every sequence with mesh tending
to zero and satisfying all the preceding domain, marked
Carath\'eodory, angle and nome assumptions,
\[
 \mathcal L(\chi_\delta\circ\gamma_\delta)
       \ \Longrightarrow\ \chi_*\mu_{\Omega,a,b,m}
\]
in the metric~\eqref{mod:curve-metric}, along the full sequence.
The physical prime-end curve law $\mu_{\Omega,a,b,m}$ is independent
of the lattices, admissible domain approximations and nome sequence.

For $m>0$, $\mu_{\Omega,a,b,m}$ is the unique law in
Definition~\ref{def:law}. For $m<0$, it is the reversal of the
positive-mass law in $(\Omega,b,a)$ with mass $|m|$, with primal and
dual boundary labels interchanged. In either case we call it
thermal massive $\mathrm{SLE}_{16/3}$.
\end{theorem}

The remainder of the paper proves the theorem. Until the final
duality argument, $m>0$.

\subsection{The scalar boundary value problem}
Our Laplacian is the ordinary Euclidean $\Delta$. Green functions
are for $-\Delta$, so their logarithmic singularity has coefficient
$1/(2\pi)$. We write $D=2\ii\partial_z$. Consequently
\begin{equation}\label{mod:square-gradient}
 \dd h=\Ima(f^2\dd z)\quad\Longleftrightarrow\quad Dh=f^2,
 \qquad |\grad h|=|f|^2.
\end{equation}

For a bounded simply connected marked domain $(O,c,d)$, use a
conformal map taking $c,d$ to $0,\infty$. Its negative real bank
is the primal-wired side, and its positive real bank is the
dual-wired side. Consider
\begin{equation}\label{mod:bvp}
 \Delta h=-4m|\grad h|\quad\hbox{in }O,\qquad
 h=1\text{ on the negative side},\quad
 h=0\text{ on the positive side}.
\end{equation}
The solution class consists of bounded functions classical in the
interior, with the specified prime-end limits off the two marks.

\begin{proposition}\label{mod:bvp-unique}
Equation~\eqref{mod:bvp} has exactly one solution in this class.
It satisfies $0\le h\le1$. On interior compact sets, $h$ and its
first derivatives depend continuously on the marked domain in the
Carath\'eodory topology, when the domains lie in a fixed bounded
physical region.
\end{proposition}

\begin{proof}
Park's two-point observable supplies a massive-holomorphic $f$ and
its square integral $h$, with the asserted arc values, bounds and
interior regularity; see \cite[Theorem~1.1, Definition~5.2,
Proposition~5.8]{Park}. Equation~\eqref{mod:square-gradient} gives
the displayed scalar equation with precisely the coefficient $4m$.
The source states uniqueness in its observable and square-integral
class. We verify that the bounded scalar class used here is also
unique.

If $h_1,h_2$ are two bounded classical solutions and $q=h_1-h_2$,
define
\[
 \beta=4m\frac{\grad h_1+\grad h_2}
                    {|\grad h_1|+|\grad h_2|},
\]
with value zero when both gradients vanish. The identity
$(|p|-|p'|)(|p|+|p'|)=(p+p')\cdot(p-p')$ shows that
\[
 \Delta q+\beta\cdot\grad q=0,\qquad |\beta|\le4m.
\]
Extend $\beta$ boundedly and measurably outside $O$ and use the
diffusion $\dd X=\sqrt2\,\dd B+\beta(X)\dd t$, killed on exit
from $O$. Bounded-drift Girsanov constructs this diffusion and makes
its law equivalent to Brownian law on every finite time interval.
Its exit time is almost surely finite. Indeed, in a fixed time block
a Brownian coordinate increment larger than the diameter of $O$
plus the maximal drift displacement forces exit. This event has a
positive probability uniform over the starting point and the past;
iteration gives a geometric bound on the number of such blocks.

Brownian exit from a simply connected domain has a prime-end limit,
and its harmonic measure gives mass zero to each specified prime
end. Finite-horizon equivalence and then a countable union over
horizons transfer these assertions to $X$. In particular its exit
prime end is almost surely different from $c,d$. It\^o's formula on
compact interior exhaustions makes $q(X)$ a bounded stopped
martingale. At exit its limit is zero by the prescribed boundary
limits. Bounded convergence and optional stopping give $q=0$.

For continuity, take any converging sequence of bounded marked
domains. Approximate each domain by a rhombic domain sufficiently
finely that its discrete square integral and reconstructed squared
fermion approximate $h$ and $Dh$ on the first finitely many compact
sets. The mesh can also be chosen so that this diagonal
sequence converges to the limiting marked domain. Park's convergence
theorem applies to that sequence. Uniqueness just proved identifies
every subsequential limit with the scalar solution in the limiting
domain. Interior estimates for the massive equation give convergence
of first derivatives on smaller compact sets. This proves the stated
continuity and does not impose a factorization condition on the
scalar solution class.
\end{proof}

\subsection{A continuum definition}
Choose a conformal chart $\phi:\Omega\to\Hh$ with
$\phi(a)=0$ and $\phi(b)=\infty$. For a continuous real driving
function $W$ with $W_0=0$, let
\begin{equation}\label{mod:loewner}
 \partial_tg_t(z)=\frac{2}{g_t(z)-W_t},\qquad g_0(z)=z,
\end{equation}
be the chordal Loewner maps, normalized at infinity, and let $O_t$
be the physical active domain, the inverse under $\phi$ of the
unbounded Loewner component. Thus $g_t$ acts in the initial
half-plane; $g_t\circ\phi$ is the physical uniformization.
Put
\begin{equation}\label{mod:chart}
 \begin{aligned}
 Z_t(w)&=g_t(\phi(w))-W_t,&\qquad
 \psi_t(v)&=\phi^{-1}(g_t^{-1}(W_t+v)),\\
 \mu_t(v)&=m|\psi_t'(v)|,&
 P(v)&=\Ima\frac{-1}{v}.
 \end{aligned}
\end{equation}
Let $h_t$ be the solution of Equation~\eqref{mod:bvp} in $O_t$,
with the two banks read in the chart $Z_t$. Define the nonnegative
extended integral
\begin{equation}\label{mod:C}
 C_t=\frac4\pi\int_{\Hh}P(v)\mu_t(v)
             |\grad(h_t\circ\psi_t)(v)|\,\dd A(v).
\end{equation}
An arbitrary fixed dummy value may be assigned when this integral
is infinite or undefined; admissibility below separately requires
that the actual integral be finite at almost every time.

\begin{definition}\label{def:law}
A positive-mass interface law in $(\Omega,a,b)$ is a law of an
oriented continuous curve in the prime-end compactification, from
$a$ to $b$, with the following properties. In the chart $\phi$ it
is the Loewner trace for the entire evolution
\eqref{mod:loewner}, parametrized by half-plane capacity $2t$,
$0\le t<\infty$, and completed continuously at $b$ in disc
coordinates. Its driving function is a weak solution of
\begin{equation}\label{mod:SDE}
 \dd W_t=\sqrt{\frac{16}{3}}\,\dd B_t
                  +\frac{2\pi}{3}C_t\,\dd t,
 \qquad W_0=0,
\end{equation}
where $B$ is Brownian motion in a filtration to which the driver is
adapted. The actual integral in Equation~\eqref{mod:C} is finite
almost everywhere and, for every finite $T$,
\begin{equation}\label{mod:finite-energy}
 \int_0^T C_t^2\,\dd t<\infty\qquad\hbox{almost surely}.
\end{equation}
The trace requirement includes the full parametrized evolution;
one cannot append curve pieces inside already swallowed hulls.
\end{definition}

This is a continuum prescription involving no lattice or selected
subsequence. Existence will follow from compactness and identification,
and uniqueness in law is proved in Section~\ref{sec:uniqueness}.
Pathwise uniqueness is neither asserted nor used.

\begin{remark}\label{mod:chart-independence}
The choice of $\phi$ is immaterial. Replacing it by $\lambda\phi$,
$\lambda>0$, replaces $t,W_t,v$ by
$\lambda^2t,\lambda W_t,\lambda v$. Directly from
Equation~\eqref{mod:C}, the new coefficient at time $\lambda^2t$
is $C_t/\lambda$. Brownian scaling therefore preserves
Equation~\eqref{mod:SDE} and its finite-energy condition. The initial
disc normalization similarly provides coordinates for the curve
topology, rather than extra data for the law.
\end{remark}
